% ===============================================================
%  PART III -- Existing Solutions to the Forced-Writeback Problem
%  Chapters 11--14. Uses the shared preamble and refs.bib.
% ===============================================================

\part{Existing Solutions to the Forced-Writeback Problem}
\label{part:solutions}

\begin{center}
\adjustbox{max width=\linewidth}{%
\begin{forest}
  part map,
  [Part III
    [Ch.~11: Owner-Based Solutions
      [Berkeley ownership] [MOESI and the O state] [Limits of ownership]]
    [Ch.~12: Other Extensions
      [Write-update: Dragon {\&} Firefly] [MESIF: the F state]]
    [Ch.~13: Scaling Approaches
      [Directories: DASH] [Token coherence]]
    [Ch.~14: Comparative Analysis
      [Mechanisms] [Traffic model] [Published evaluations] [Research gap]]
  ]
\end{forest}}
\end{center}

\begin{keybox}[What Part III does]
Part~III reviews published approaches against the problem of Chapter~10. For each
one we ask what it does when a block held in \stM\ is read by another cache, what
that costs, and whether it meets the requirements of \cref{fig:req-tree}. The
approaches fall into three classes. Direct solutions change the block states so
that modified data can be shared without updating memory (Chapter~11).
Alternative strategies change how shared data is propagated or forwarded, and
affect the forced writeback as a side effect (Chapter~12). Organisations change
how requests travel and are ordered, and leave the writeback to the state
protocol that runs on top of them (Chapter~13). Chapter~14 compares all three.
\end{keybox}

% ===============================================================
%  Chapter 11 -- Owner-Based Solutions: MOESI and the Berkeley Protocol
% ===============================================================
\chapter{Owner-Based Solutions: MOESI and the Berkeley Protocol}
\label{ch:owner}
\index{MOESI}\index{Berkeley protocol}\index{owner}%

\begin{figure}[htbp]
\centering
\adjustbox{max width=\linewidth}{%
\begin{forest}
  chapter map,
  [Chapter 11
    [Motivation [Supplier vs owner vs memory] [Clean-S is a choice]]
    [Berkeley protocol [Shared-Dirty owner] [Invalidate before write]]
    [MOESI [Three attributes] [M to O] [Futurebus signals]]
    [Comparison {\&} limits [Table A] [Deferred not removed]]
  ]
\end{forest}}
\caption[]{Structure of Chapter~11. These are the only
approaches in Part~III designed to let modified data be shared without first
updating memory.}
\label{fig:ch11-map}
\end{figure}

% ---------------------------------------------------------------
\section{Motivation for Owner-Based Protocols}
\label{sec:own-motivation}
\sectionmap{[Motivation [Three roles for one block] [MESI merges them on M to S]
  [Ownership separates responsibility from memory]]}
\index{MESI}\index{MOESI}\index{Berkeley protocol}\index{coherence axioms}\index{owner}%

A cached block involves three roles: the supplier, which gives the current data
to a requester; the responsible party, which must make sure the current data
reaches memory in the end; and memory, which may or may not hold the current
value. Under MESI the responsible party is the one cache holding the block in
\stM\ or, if there is none, memory. A block in \stS\ is therefore valid only if
memory is current (the clean-shared rule of \cref{sec:dirty-root}), and the
transition $\stM\to\stS$ has to hand responsibility to memory by writing the block
back. As an invariant, let $\mathcal{O}(b)$ be the set of caches responsible for
block $b$. Every protocol in this chapter keeps
\begin{equation}
  |\mathcal{O}(b)| \le 1
  \qquad\text{and}\qquad
  \mathcal{O}(b) = \varnothing \;\Rightarrow\; \text{memory holds the current value of } b.
  \label{eq:owner-invariant}
\end{equation}
MESI also requires that no other valid copy exist while $\mathcal{O}(b)$ is
non-empty, and it is this extra rule, not \cref{eq:owner-invariant}, that forces
the writeback. Neither coherence axiom (\cref{sec:coh-meaning}) needs memory to be
current while some cache is still responsible for the block. An owner-based
protocol drops only the extra rule: one cache, the \gls{owner}, stays responsible
for a modified block while other caches hold read-only copies, and memory is
updated when ownership ends (\cref{fig:roles}).

\begin{figure}[htbp]
\centering
\begin{bigfig}[0.88]
\begin{tikzpicture}[font=\sffamily\scriptsize,
    ag/.style={draw, thick, rounded corners=2pt, minimum width=1.9cm,
               minimum height=0.65cm, align=center, font=\sffamily\small},
    resp/.style={circle, fill=moesiOrange, text=white, font=\sffamily\bfseries\tiny,
                 inner sep=1.2pt},
    m/.style={-{Stealth}, thick}]
  \node[font=\sffamily\bfseries\small] at (2.2,3.0) {(a) MESI: responsibility moves to memory};
  \node[ag, fill=mesiBlue!15, draw=mesiBlue] (A) at (0,2.0) {Cache A\\\stM$\,\to\,$\stS};
  \node[ag, fill=mesiBlue!15, draw=mesiBlue] (B) at (4.4,2.0) {Cache B\\\stI$\,\to\,$\stS};
  \node[ag, fill=memGreen!15, draw=memGreen] (M) at (2.2,0.0) {Memory\\current};
  \draw[m, mesiBlue] (A) -- node[above]{data} (B);
  \draw[m, alertRed, very thick] (A) -- node[left, text=alertRed]{write back} (M);
  \node[resp] at ($(M.north east)+(0.05,0.05)$) {R};
  \begin{scope}[xshift=7.6cm]
  \node[font=\sffamily\bfseries\small] at (2.2,3.0) {(b) Ownership: responsibility stays in A};
  \node[ag, fill=moesiOrange!15, draw=moesiOrange] (A2) at (0,2.0) {Cache A (owner)\\\stM$\,\to\,$\stO};
  \node[ag, fill=moesiOrange!15, draw=moesiOrange] (B2) at (4.4,2.0) {Cache B\\\stI$\,\to\,$\stS};
  \node[ag, fill=coreGrey!10, draw=coreGrey] (M2) at (2.2,0.0) {Memory\\stale};
  \draw[m, moesiOrange] (A2) -- node[above]{data} (B2);
  \draw[dashed, coreGrey] (A2) -- (M2);
  \node[text=coreGrey, align=right, anchor=east] at (0.9,0.45) {written back only\\when A evicts\\the block};
  \node[resp] at ($(A2.north east)+(0.05,0.05)$) {R};
  \end{scope}
\end{tikzpicture}
\end{bigfig}
\caption[Responsibility for a dirty block after a remote read]{Where responsibility
(R) for a modified block sits after another cache reads it. Under MESI it passes to
memory through a writeback; under ownership it stays with the supplier and memory
stays stale.}
\label{fig:roles}
\end{figure}

The idea is older than the MOESI name. Archibald and Baer, describing the
protocols of the early 1980s, say the Berkeley protocol ``uses the idea of
ownership'', and the owner they define is the cache that has the block in
\textsc{Dirty} or \textsc{Shared-Dirty}~\cite[pp.~277--278]{archibald1986evaluation}.

% ---------------------------------------------------------------
\section{The Berkeley Ownership Protocol}
\label{sec:own-berkeley}
\sectionmap{[Berkeley protocol [Designed with its controller chip] [Four states]
  [Owner supplies and keeps ownership] [Memory updated on replacement]
  [Invalidation before every non-owner write]]}
\index{MESI}\index{MSI}\index{Exclusive state}\index{forced writeback}\index{snooping protocol}\index{Berkeley protocol}\index{coherence axioms}\index{memory consistency}\index{owner}%

Katz, Eggers, Wood, Perkins and Sheldon designed the Berkeley protocol and
published it as a technical report~\cite{katz1984implementing} and at
ISCA~1985~\cite{katz1985implementing}. They call it an ``ownership-based
multiprocessor cache consistency protocol''. The report presents it together
with a cache controller meant to fit on one VLSI chip, spends much of its length
on making the controller's concurrent operations safe, and takes that design down
to a CMOS layout~\cite{katz1984implementing}. The state set is therefore one a
snooping controller of the early 1980s could actually implement. We could not read
the original report in full (see the note at the end of this chapter), so the
states below follow Archibald and Baer's
description~\cite[pp.~277--278]{archibald1986evaluation} (\cref{tab:berkeley-states}).
There is no counterpart of \stE: a block read when no other cache holds it enters
\textsc{Valid}, as in MSI.

\begin{table}[htbp]
\centering
\caption{States of the Berkeley protocol as described by Archibald and
Baer~\cite[pp.~277--278]{archibald1986evaluation}.}
\label{tab:berkeley-states}
\small
\begin{tabular}{@{}l L{4.2cm} c c l@{}}
\toprule
State & Meaning & Owner? & Other copies? & Nearest MOESI state \\
\midrule
\textsc{Invalid}      & no valid copy & --- & --- & \stI \\
\textsc{Valid}        & read-only copy; not responsible for memory & no & possibly & \stS \\
\textsc{Shared-Dirty} & modified; owned by this cache & yes & possibly & \stO \\
\textsc{Dirty}        & modified; owned by this cache, sole copy & yes & no & \stM \\
\bottomrule
\end{tabular}
\end{table}

On a miss to a block that another cache holds in \textsc{Dirty} or
\textsc{Shared-Dirty}, the owning cache ``must supply the block contents directly
to the other cache''. The owner then holds the block in \textsc{Shared-Dirty} and
the requester in \textsc{Valid}~\cite[p.~278]{archibald1986evaluation}. Memory is
not written, because a modified block is not written back on becoming
shared~\cite[p.~277]{archibald1986evaluation}, so responsibility stays with the
supplier as in \cref{fig:roles}(b). Memory catches up only on replacement: ``A
block in either state \textsc{Shared-Dirty} or \textsc{Dirty} must be written back
to main memory if it is selected for
replacement''~\cite[p.~277]{archibald1986evaluation}.

A cache becomes the owner by writing. A write hit to a block in \textsc{Valid} or
\textsc{Shared-Dirty} requires that ``an invalidation signal must be sent on the
bus before the write is allowed to proceed''. Every other copy is invalidated and
the writer's block becomes \textsc{Dirty}~\cite[p.~278]{archibald1986evaluation}.
If another cache was the owner, its duty ends with the invalidation, since the new
write replaces its data, and no writeback is needed.

On the transaction of \cref{fig:dirty-transaction} the Berkeley protocol removes
the forced writeback completely: one request and one data transfer replace MESI's
request, data transfer and writeback. The cost shows up elsewhere. With no \stE,
the first write to any block fetched from memory needs a bus invalidation even
when no other cache holds the block, which is the inefficiency MESI was designed
to remove (\cref{sec:fam-mesi}). Against \cref{fig:req-tree}, Berkeley keeps both
axioms and removes the forced writeback but loses MESI's silent upgrade, so it is
a ``weaker candidate''.

% ---------------------------------------------------------------
\section{MOESI: The Owned State}
\label{sec:own-moesi}
\sectionmap{[MOESI [A class of compatible protocols] [Three attributes per line]
  [Intervention: O and M supply data] [Memory updated on replacement]]}
\index{MESI}\index{MOESI}\index{Owned state}\index{Exclusive state}\index{forced writeback}\index{Berkeley protocol}\index{Dragon protocol}\index{Firefly protocol}\index{Illinois protocol}\index{Futurebus}\index{owner}%

Sweazey and Smith defined a whole class of protocols. As
long as every cache ``uses a protocol from the class we define, the overall memory
system state will remain consistent''~\cite[p.~415]{sweazey1986moesi}, and they
showed how the IEEE Futurebus could support any member of the class. Each cache
may pick, at any moment, any action that some protocol in the class
allows~\cite[abstract]{sweazey1986moesi}. They count Berkeley and Dragon as members
and show how Illinois, Firefly and Write-Once can be adapted to
fit~\cite[abstract]{sweazey1986moesi}. ``MOESI'' therefore names a state model and
a set of allowed actions. Two implementations with the same five states may still
answer the same request differently, so they have to be compared by what they do
(\cref{sec:own-limits}).

Every line has three attributes, validity, exclusiveness and ownership, and the
five legal combinations, which ``can be specified with three
bits''~\cite[p.~416]{sweazey1986moesi}, are the five states: \stM\ (exclusive,
owned), \stO\ (shareable, owned), \stE\ (exclusive, unowned, unmodified), \stS\
(shareable, unowned) and \stI. The owned states \stM\ and \stO\ are
\emph{intervenient}: a cache holding a line in either one ``when it sees a read
miss\,\ldots\ must supply the data''~\cite[pp.~418--419]{sweazey1986moesi}. The
owner must either make sure memory is eventually updated, pass ownership to
another cache, or answer requests in place of
memory~\cite[p.~417]{sweazey1986moesi}.

\Cref{fig:futurebus-read} shows the remote read on Futurebus. The requester issues
a read. In the address cycle the owner asserts \textsc{CH} (cache hit: it has the
data and will keep it) and \textsc{DI} (data intervention: it owns the data and
takes over the response from memory)~\cite[pp.~417--418]{sweazey1986moesi}. The
owner then sends the line and moves to \stO, and the requester loads it in \stS.
\Cref{tab:o-state-questions} answers, from the source, the questions this
transition raises.

\begin{figure}[htbp]
\centering
\begin{bigfig}[0.85]
\begin{tikzpicture}[font=\sffamily\scriptsize,
    ag/.style={draw, thick, rounded corners=2pt, minimum width=2.0cm,
               minimum height=0.6cm, align=center, font=\sffamily\small},
    m/.style={-{Stealth}, thick}, sg/.style={font=\sffamily\bfseries\scriptsize}]
  \node[ag, fill=moesiOrange!15, draw=moesiOrange] (B) at (0,0) {Requester};
  \node[ag, fill=moesiOrange!15, draw=moesiOrange] (A) at (5.0,0) {Owner (\stM)};
  \node[ag, fill=memGreen!15, draw=memGreen] (Mm) at (10.0,0) {Memory};
  \foreach \x in {0,5,10} \draw[dashed, coreGrey!60] (\x,-0.35) -- (\x,-4.0);
  \draw[m, coreGrey] (0,-0.9) -- node[above]{read (address cycle)} (5,-0.9);
  \draw[coreGrey, thick] (5,-0.9) -- (10,-0.9);
  \node[sg, text=moesiOrange, anchor=west] at (5.15,-1.45) {asserts CH: keeps a copy};
  \node[sg, text=moesiOrange, anchor=west] at (5.15,-1.9) {asserts DI: memory preempted};
  \draw[m, moesiOrange, very thick] (5,-2.6) -- node[above]{data (cache-to-cache)} (0,-2.6);
  \node[text=coreGrey, align=left, anchor=west] at (10.15,-2.6) {no response,\\no update};
  \node[font=\sffamily\bfseries\small, text=moesiOrange] at (0,-3.6) {\stS};
  \node[font=\sffamily\bfseries\small, text=moesiOrange] at (5,-3.6) {\stO};
\end{tikzpicture}
\end{bigfig}
\caption[The \stM$\to$\stO\ transaction on Futurebus]{A remote read of a modified
line in a MOESI-class protocol, with the Futurebus signals of Sweazey and
Smith~\cite[pp.~417--418]{sweazey1986moesi}. The owner keeps the line and the
responsibility; memory neither responds nor is updated.}
\label{fig:futurebus-read}
\end{figure}

\begin{table}[htbp]
\centering
\caption{The \stO\ state, transaction by transaction.}
\label{tab:o-state-questions}
\small
\begin{tabular}{@{}L{4.0cm} L{8.9cm}@{}}
\toprule
Question & Answer and basis \\
\midrule
Why does the owner send the data? & Only the owner has the current value, and
  intervenient caches must supply it on a read miss~\cite[pp.~418--419]{sweazey1986moesi}. \\
Why \stO\ and not \stS? & \stS\ is unowned; moving to it would leave nobody
  responsible while memory is stale, breaking \cref{eq:owner-invariant}. \\
Why may memory stay stale? & The owner remains responsible and must update memory,
  pass ownership or answer for memory~\cite[p.~417]{sweazey1986moesi}. \\
Why can \stS\ coexist with \stO? & \stS\ copies are read-only and equal to the
  owner's; no copy can be written without first invalidating the others, so the
  order axiom holds. \\
Why is the owner still responsible? & Ownership is one of the three line
  attributes and is not given up by supplying data~\cite[p.~416]{sweazey1986moesi}. \\
What happens on eviction? & Copy-back: changes reach memory ``only when the line is
  replaced''~\cite[p.~416]{sweazey1986moesi}. \\
How are writes handled? & A write needs exclusiveness; the writer invalidates other
  copies (Futurebus \textsc{IM}, intent to modify) and holds the line in
  \stM~\cite[pp.~417--418]{sweazey1986moesi}. The old owner's data is superseded and
  is not written back. \\
\bottomrule
\end{tabular}
\end{table}

MESI and MOESI therefore differ in when memory is written; both are coherent. For a
producer and one consumer exchanging a block for $k$ rounds before it is evicted,
the memory writes over the block's lifetime are
\begin{equation}
  W_{\mathrm{MESI}}(k) = k,
  \qquad
  W_{\mathrm{MOESI}}(k) = 1,
  \label{eq:wb-timing}
\end{equation}
because MESI writes back at every \stM$\to$\stS\ transition and then evicts a clean
block, while MOESI writes back once, when the owner evicts. In the notation of
\cref{eq:fwb-fraction}, MOESI has $B_{\mathrm{fwb}}=0$ by construction. Its single
deferred writeback is a different quantity, $B_{\mathrm{dwb}}$, paid at eviction.
\Cref{eq:wb-timing} is our own count under the assumptions of Part~II.

In hardware, the \stO\ state needs three bits per line instead of
two~\cite[p.~416]{sweazey1986moesi}, and the bus must let a cache take over from
memory (\textsc{DI}) and announce that it keeps a copy (\textsc{CH}). Sweazey and
Smith also point out that Futurebus ``has no mechanism by which a transfer from one
cache to another can also update main memory''. Protocols that need that update,
including Write-Once and Firefly, have to abort the transaction with \textsc{BS}
(busy), write the line back and start again~\cite[p.~418]{sweazey1986moesi}. On
such a bus a MESI-style forced writeback costs an aborted and restarted
transaction, which is more than one extra transfer.

% ---------------------------------------------------------------
\section{Comparison and Limitations of Owner-Based Approaches}
\label{sec:own-compare}
\label{sec:own-limits}
\sectionmap{[Comparison and limits [Same ownership idea] [Berkeley lacks E]
  [Deferred not eliminated] [Implementations differ] [Benefit depends on workload]]}
\index{MESI}\index{MOESI}\index{MESIF}\index{Owned state}\index{Exclusive state}\index{forced writeback}\index{directory protocol}\index{Berkeley protocol}\index{Illinois protocol}\index{Futurebus}\index{gem5}\index{owner}%

Berkeley and MOESI use the same mechanism: on a remote read the owner supplies the
data, keeps ownership in a shared-dirty state, and leaves memory stale until
replacement. They differ in three ways. Berkeley has no \stE, so the first write
to a block fetched from memory costs a bus invalidation. Berkeley is one protocol
built for one controller, while MOESI is a class whose members may choose
different allowed actions. And the Futurebus version defines the bus signals
(\textsc{CH}, \textsc{DI}) that assert ownership. \Cref{tab:ownership-protocols}
adds the other published protocols with a shared-dirty state.

\begin{table}[htbp]
\centering
\caption{Table A: ownership-based protocols and shared-dirty variants.}
\label{tab:ownership-protocols}
\scriptsize
\begin{tabular}{@{}L{1.55cm} L{2.0cm} L{2.15cm} L{1.9cm} L{2.4cm} L{2.4cm}@{}}
\toprule
Protocol & State representation & Owner on remote read & Memory updated &
  Additional messages / signals & Principal limitation \\
\midrule
Berkeley~\cite{archibald1986evaluation} & \textsc{I, Valid, SD, Dirty};
  2 bits & supplies; \textsc{Dirty}/\textsc{SD}$\to$\textsc{SD}; requester
  \textsc{Valid} & on replacement of \textsc{SD} or \textsc{Dirty} &
  invalidation before any write to \textsc{Valid} or \textsc{SD} & no \stE:
  invalidation on first write even if unshared \\
MOESI class~\cite{sweazey1986moesi} & M, O, E, S, I; 3 bits & supplies
  (intervention); $\stM\to\stO$ & on replacement, or ownership passed &
  \textsc{CH}, \textsc{DI} bus signals & members may choose different
  actions; Futurebus cannot update memory during a transfer \\
Dragon~\cite{archibald1986evaluation} & \textsc{VE, SC, SD, Dirty};
  no invalid state & supplies; $\to$\textsc{SD}; requester \textsc{SC} &
  on replacement of \textsc{SD} & a bus update on every write to a shared
  block & update traffic per write (Chapter~12) \\
MESIF variant FM~\cite{goodman2009mesif} & proposed sixth state,
  Shared-Modified & requester may hold FM and defer cleansing & deferred;
  not specified & not specified & described in one paragraph only \\
MOESI-prime~\cite{loughlin2022moesiprime} & MOESI + M$'$, O$'$; 7 states,
  3 bits & as MOESI & as MOESI & avoids redundant in-DRAM directory writes
  & added directory-cache contention \\
\bottomrule
\end{tabular}
\end{table}

Owner-based protocols remove the forced writeback, but the memory update still has
to happen, and the literature gives no grounds for treating them as always better
than MESI. Five limits stand out.

\paragraph{The writeback is deferred.} The owner still writes the block back when
it evicts it (\cref{eq:wb-timing}). The saving is the $k-1$ overwritten writebacks
of \cref{sec:dirty-wasted}. If a block is shared once and then evicted, both
protocols make the same single write at different times. Whether the later write
is cheaper depends on when eviction happens; we pose this as RQ4 (\cref{tab:rq}),
and none of the sources answers it.

\paragraph{Memory is stale while an owner exists.} Every request for the block must
be answered by the owner, which must stand in for memory when
responding~\cite[p.~417]{sweazey1986moesi}. In our analysis this puts every read
of a widely shared dirty block on one cache's port, where MESI would let later
readers be served by memory or any sharer.

\paragraph{More state and control.} MOESI needs three bits per line instead of
two~\cite[p.~416]{sweazey1986moesi} and a bus that can preempt memory. The Berkeley
design needed explicit work on controller critical sections and
interlocks~\cite{katz1984implementing}, since races appear when ownership moves
during concurrent requests.

\paragraph{Implementations differ.} A protocol called MOESI need not perform the
\stM$\to$\stO\ transition. Mallya et al.\ report that the gem5
\texttt{MOESI\_CMP\_directory} implementation they used invalidates its copy on a
forwarded read of an \stM\ block instead of sharing dirty data, and that it differs
from \texttt{MESI\_Two\_Level} in many other states and
transitions~\cite{mallya2015simulation}. A ``MOESI versus MESI'' result has to be
read with the implementation in mind.

\paragraph{The benefit depends on the workload.} Every source that compares an
owner protocol with an invalidate protocol finds the advantage conditional.
Archibald and Baer report that Berkeley ``actually surpasses the Illinois scheme
at levels of high sharing'', because on a miss to a modified block Berkeley ``does
not require updating main memory as does
Illinois''~\cite[p.~295]{archibald1986evaluation}. Fagundes and Souza find that
MOESI's advantage over MESI shrinks as the LLC grows~\cite{fagundes2024evaluation}.
Loughlin et al.\ find that replacing MESI with their owner-based MOESI-prime cuts
the peak DRAM row activation rate substantially while changing execution time by
only $-0.51\,\%$ to $+0.61\,\%$~\cite{loughlin2022moesiprime}, so fewer memory
writes need not mean a shorter run.

\begin{remarkbox}[A note on sources]
We identified the original Berkeley report~\cite{katz1984implementing} and
paper~\cite{katz1985implementing} and checked their bibliographic details, but
could not read their full text. The state names and transitions above therefore
follow Archibald and Baer~\cite{archibald1986evaluation}. Other sources say the
original paper uses the terms UnOwned, Owned Exclusively and Owned NonExclusively;
we have not been able to check this against it.
\end{remarkbox}

% ===============================================================
%  Chapter 12 -- Other Extensions: MESIF and Write-Update Protocols
% ===============================================================
\chapter{Other Extensions: MESIF and Write-Update Protocols}
\label{ch:other-ext}
\index{MESIF}%

\begin{figure}[htbp]
\centering
\adjustbox{max width=\linewidth}{%
\begin{forest}
  chapter map,
  [Chapter 12
    [Write-update protocols [Firefly: caches and memory] [Dragon: caches only]
      [Update vs invalidate traffic]]
    [MESIF [The F state] [Reads of M lines migrate]]
    [Discussion [Remote-read decision tree] [Table B]]
  ]
\end{forest}}
\caption[]{Structure of Chapter~12. Neither family was
designed to remove the forced writeback; we look at what each does to it as a side
effect.}
\label{fig:ch12-map}
\end{figure}

% ---------------------------------------------------------------
\section{Write-Update Protocols}
\label{sec:ext-update}
\sectionmap{[Write-update [Writes propagate to sharers] [Firefly: also to memory]
  [Dragon: caches only with an owner] [Trade: no misses for update traffic]]}
\index{MESI}\index{MOESI}\index{forced writeback}\index{write-back cache}\index{write-through cache}\index{Berkeley protocol}\index{Dragon protocol}\index{Firefly protocol}\index{owner}%

A \gls{update} protocol keeps every cached copy valid by sending each write to the
other caches that hold the block (\cref{sec:coh-update-invalidate}). Archibald and
Baer call this the distributed-write approach and found that in their model it
``yields the best performance in the handling of shared
data''~\cite[p.~294]{archibald1986evaluation}. Firefly and Dragon are often
grouped together, but they update memory and assign ownership differently, and
that difference decides what each does to the forced writeback.

\paragraph{Firefly.} The DEC Firefly cache has two tag bits per line,
\textsc{Dirty} and \textsc{Shared}, and the bus has a signal, \textsc{MShared},
``asserted during an operation if another cache contains the requested
location''~\cite[p.~7]{thacker1987firefly}. A write to an unshared line is handled
as in any write-back cache and ``requires no MBus traffic''. A shared line is
written through: ``the cache does write-through, and other caches that share the
datum are updated, as is main storage''~\cite[p.~7]{thacker1987firefly}. On a read
miss, when another cache holds the line, ``the caches that contain the line supply
the data, and the memory is inhibited''~\cite[p.~8]{thacker1987firefly}. If the
supplying cache holds the line dirty, Archibald and Baer say the block is supplied
and ``written to main memory at the same time'', after which every cache holding
it, the requester included, is in \textsc{Shared}~\cite[p.~280]{archibald1986evaluation}.
Firefly therefore never lets a modified block be shared. On the transaction of
\cref{fig:dirty-transaction} it makes the same memory update as MESI, and after
that every write to the shared block also goes to memory: ``the Firefly must send
distributed writes to memory while the Dragon sends them to the caches
only''~\cite[p.~295]{archibald1986evaluation}. Firefly keeps the forced writeback
and adds a memory write for each later shared write.

\paragraph{Dragon.} Dragon's states are \textsc{Valid-Exclusive},
\textsc{Shared-Clean}, \textsc{Shared-Dirty} and \textsc{Dirty}, with no invalid
state~\cite{archibald1986evaluation}. On a read miss to a block held elsewhere in
\textsc{Dirty} or \textsc{Shared-Dirty}, the holder ``supplies the data, raises the
SharedLine, and sets its block state to \textsc{Shared-Dirty}'', and the requester
takes the block as \textsc{Shared-Clean}~\cite[p.~281]{archibald1986evaluation}.
Writes to a shared block are broadcast to the other holders, and ``the single-word
write does not go to main memory''~\cite[p.~282]{archibald1986evaluation}. ``The
last cache to write it, if any, will have the block in state \textsc{Shared-Dirty}''
and is ``responsible for updating memory when the block is
replaced''~\cite[p.~281]{archibald1986evaluation}. Dragon thus has the same
ownership idea as Berkeley and MOESI, a single shared-dirty owner responsible for
memory, combined with update in place of invalidation, and on a remote read of a
modified block it defers the writeback as MOESI does. Sweazey and Smith place
Dragon in the MOESI class and note that its memory is updated ``only on a
replacement''~\cite[p.~420]{sweazey1986moesi}.

\paragraph{Update traffic against forced writebacks.} Without invalidations there
are no misses after a write, but each write to a shared block costs a bus
transaction, so the outcome depends on how many writes the producer makes between
reads. We add one assumption to the counting rules of \cref{tab:msg-classes}: a
single-word update is an 8-byte header plus an 8-byte word, $u=16$\,B. For one
producer, $C$ consumers and $w$ producer writes per round, the steady-state bytes
per round are
\begin{align}
  B_{\mathrm{MESI}}(C)            &= 160 + 80\,(C-1), \label{eq:upd-mesi}\\
  B_{\mathrm{MOESI}}(C)           &= \phantom{1}88 + 80\,(C-1), \label{eq:upd-moesi}\\
  B_{\mathrm{Dragon}}(w)          &= B_{\mathrm{Firefly}}(w) = u\,w = 16\,w. \label{eq:upd-dragon}
\end{align}
\Cref{eq:upd-mesi} is \cref{eq:pc-model}. In \cref{eq:upd-moesi} the owner serves
the consumer's read without a writeback ($8+72$\,B) and the producer's write is an
invalidating upgrade (8\,B). In \cref{eq:upd-dragon} the consumers' reads all hit
and one bus update reaches every sharer, so the cost does not depend on $C$.
Firefly sends the same bus bytes, but each of its $w$ updates also writes memory.
These are our analytical counts. \Cref{fig:update-vs-invalidate} explains how two
observations that seem to clash can both hold. Archibald and Baer, whose synthetic
workload made single-word references to shared blocks, found distributed write
best for shared data~\cite[p.~294]{archibald1986evaluation}, and Chapter~4 gave the
usual argument that invalidation wins when a core writes a block many times before
others read it. The two results sit on opposite sides of the crossover.

\begin{figure}[htbp]
\centering
\begin{tikzpicture}
  \begin{axis}[width=10.5cm, height=5.2cm,
      xlabel={Writes by the producer per round, $w$},
      ylabel={Bus bytes per round ($C=1$)},
      xmin=1, xmax=16, ymin=0, ymax=270, xtick={1,2,4,6,8,10,12,14,16},
      legend style={at={(0.02,0.98)}, anchor=north west, font=\scriptsize},
      grid=major, grid style={coreGrey!15},
      label style={font=\small}, tick label style={font=\scriptsize}]
    \addplot[very thick, mesiBlue, domain=1:16] {160};
    \addlegendentry{MESI (invalidate, forced write back)}
    \addplot[very thick, moesiOrange, domain=1:16] {88};
    \addlegendentry{MOESI / Berkeley (invalidate, owner)}
    \addplot[very thick, tealC, domain=1:16] {16*x};
    \addlegendentry{Dragon / Firefly (update)}
    \draw[dashed, coreGrey] (axis cs:5.5,0) -- (axis cs:5.5,88);
    \draw[dashed, coreGrey] (axis cs:10,0) -- (axis cs:10,160);
    \node[font=\scriptsize, text=coreGrey, anchor=south west] at (axis cs:5.5,4) {5.5};
    \node[font=\scriptsize, text=coreGrey, anchor=south west] at (axis cs:10,4) {10};
  \end{axis}
\end{tikzpicture}
\caption[Update and invalidate traffic per producer--consumer round]{Bus bytes per
round for one producer and one consumer
(\cref{eq:upd-mesi,eq:upd-moesi,eq:upd-dragon}). Update protocols move fewer bytes
than MESI below ten writes per round and fewer than MOESI below 5.5; above those
points invalidation is cheaper. Firefly also writes memory once per update. Our
analytical model.}
\label{fig:update-vs-invalidate}
\end{figure}

% ---------------------------------------------------------------
\section{MESIF: The Forward State}
\label{sec:ext-mesif}
\sectionmap{[MESIF [Point-to-point interconnect] [F: one clean forwarder]
  [A read of M migrates the line] [F is not O]]}
\index{MOESI}\index{MESIF}\index{forced writeback}\index{snooping protocol}\index{directory protocol}\index{sharing patterns!producer--consumer}%

Goodman and Hum describe \gls{mesif}, finished in 2001 and used in Intel's
QuickPath Interconnect, as a way to get the latency of snooping on a point-to-point
network~\cite[\S1]{goodman2009mesif}. Their reference point is a directory
protocol, which needs ``a minimum of three serial communications to complete a
memory transaction'' where snooping needs two~\cite[\S1]{goodman2009mesif}. When
several caches hold a clean line in \stS, none of them is named to answer, so the
home memory has to serve the request. The Forward state is ``a state assigned to a
single node among those sharing a cache line'', and that node is ``responsible for
providing its content when a request is received''~\cite[\S2.1]{goodman2009mesif}.
When the F holder answers a read, the requester gets the line in F and the old
holder drops to \stS~\cite[\S2.1.1, case ii]{goodman2009mesif}, so the role moves to
the latest reader. The authors compare F to an owned state, ``but coherent with
memory and useful for read requests rather than write
requests''~\cite[\S2.1]{goodman2009mesif}.

The report treats a read of a line that another node holds in \stM\ as its hardest
case. In the basic protocol ``the line is transmitted to the requesting node as if
the request had been a Read-for-Ownership'': the holder gives up its copy and the
requester gets the line exclusively~\cite[\S2.1.1, case iv]{goodman2009mesif}. No
memory update is described for this case. A variant lets the holder keep a copy in
\stS; the requester must then either ``cleanse'' the line, meaning write it back,
or keep it in a sixth state, Shared-Modified (FM), which postpones the write. The
report says nothing more about FM~\cite[\S2.1.1, case iv]{goodman2009mesif}.

The basic behaviour avoids the forced writeback by another route than MOESI:
modified data is never shared because the line migrates. In a producer--consumer
exchange this costs something MOESI does not pay. By our analysis the producer
loses its copy at every consumer read, so its next write is a miss that must fetch
the data (an 8-byte request and a 72-byte transfer) instead of an 8-byte upgrade,
and with $C$ consumers the line migrates on every read. The paper's projected gain
of 6 to 11\,\% for four nodes comes from an in-house mean-value-analysis model whose
algorithms, the authors say, ``cannot be divulged''~\cite[\S3]{goodman2009mesif}.
For the hitI and hitS cases ``MESIF offers no benefits over the 3-hop
protocols''~\cite[\S3]{goodman2009mesif}. The protocol adds snoop stalls while a
forwarding node waits for an acknowledgement, needs the home node to identify every
party to a conflict ``even in the case of read-only (shared) data'', and uses three
virtual channels per direction to avoid deadlock~\cite[\S2.1--2.3]{goodman2009mesif}.
\Cref{tab:f-vs-o} sets F against O.

\begin{table}[htbp]
\centering
\caption{The Forward state and the Owned state compared.}
\label{tab:f-vs-o}
\small
\begin{tabular}{@{}L{3.6cm} L{4.3cm} L{4.3cm}@{}}
\toprule
 & F (MESIF) & O (MOESI) \\
\midrule
Data relative to memory & clean: ``coherent with memory''~\cite[\S2.1]{goodman2009mesif}
  & dirty: memory may be stale~\cite[p.~416]{sweazey1986moesi} \\
Responsible for write back & no & yes \\
Purpose & choose one cache to answer reads of clean shared data &
  let modified data be shared without updating memory \\
Effect on the forced writeback & none; the \stM\ case is handled by
  migration & removes it; deferred to eviction \\
Interconnect assumed & point-to-point with a home node & designed for a
  shared bus (Futurebus) \\
\bottomrule
\end{tabular}
\end{table}

% ---------------------------------------------------------------
\section{Limitations and Discussion}
\label{sec:ext-discussion}
\sectionmap{[Discussion [Five answers to one transaction] [Update: trades misses for writes]
  [Forwarding: a latency fix] [Table B]]}
\index{Dragon protocol}\index{Firefly protocol}%

\Cref{fig:remote-read-tree} sorts the protocols of Chapters~11 and~12 by what they
do on the one transaction this report studies. The sources show five different
behaviours, and only one of them both shares the data and leaves the holder
responsible for it. Write-update and forwarding target other costs than the one in
Chapter~10. Update protocols remove the misses caused by invalidation and pay with
a bus transaction per shared write (\cref{fig:update-vs-invalidate}); Dragon also
defers memory updates through ownership and Firefly does not. Forwarding reduces
the latency of serving clean shared data on a point-to-point network, and handles
modified data by migration at the cost of the producer's copy.
\Cref{tab:alt-strategies} summarises the three.

\begin{figure}[htbp]
\centering
\adjustbox{max width=0.9\linewidth}{%
\begin{forest}
  decision tree,
  [Remote read of a modified block:\\does the holder supply the data?
    [{Synapse: holder sends NAK\\writes back; requester\\retries from memory}, el=no]
    [Does the holder\\keep a copy?, el=yes
      [{MESIF (basic): line\\migrates as a\\read-for-ownership}, el=no]
      [Is memory updated\\at this point?, el=yes
        [{Illinois / MESI\\Write-once, Firefly:\\forced write back}, el=yes]
        [{Berkeley, MOESI\\Dragon: holder stays\\responsible owner}, el=no]]]]
\end{forest}}
\caption[Decision tree: remote read of a modified block]{Five published behaviours
on the transaction studied here. Synapse, Write-once, Illinois, Firefly, Berkeley
and Dragon from Archibald and Baer~\cite[pp.~276--281]{archibald1986evaluation};
MESIF from Goodman and Hum~\cite[\S2.1.1]{goodman2009mesif}; MOESI from Sweazey and
Smith~\cite[pp.~416--419]{sweazey1986moesi}.}
\label{fig:remote-read-tree}
\end{figure}

\begin{table}[htbp]
\centering
\caption{Table B: alternative coherence strategies and the forced writeback.}
\label{tab:alt-strategies}
\scriptsize
\begin{tabular}{@{}L{1.3cm} L{1.7cm} L{2.3cm} L{2.3cm} L{2.0cm} L{2.6cm}@{}}
\toprule
Protocol & Strategy & State or forwarding mechanism & Shared modified data &
  Memory updated & Relevance to the forced writeback \\
\midrule
Firefly & update, write-through when shared & \textsc{Dirty} and
  \textsc{Shared} tag bits; \textsc{MShared} line~\cite{thacker1987firefly} &
  not allowed; dirty block written to memory when supplied~\cite{archibald1986evaluation}
  & on every shared write and on eviction & retains it, and adds a memory
  write per shared write \\
Dragon & update, caches only & \textsc{VE, SC, SD, Dirty};
  \textsc{SD} owner~\cite{archibald1986evaluation} & allowed (\textsc{SD})
  & on replacement of \textsc{SD}~\cite{archibald1986evaluation} & removes it
  (defers), but adds an update per shared write \\
MESIF & invalidate, clean forwarding & F: one forwarder among clean
  copies~\cite{goodman2009mesif} & not in the basic protocol; FM variant
  proposed & not stated for the basic \stM\ case & addresses a related
  problem; avoids it by migration in the basic protocol \\
\bottomrule
\end{tabular}
\end{table}

% ===============================================================
%  Chapter 13 -- Scaling Approaches: Directories and Token Coherence
% ===============================================================
\chapter{Scaling Approaches: Directories and Token Coherence}
\label{ch:scaling-approaches}

\begin{figure}[htbp]
\centering
\adjustbox{max width=\linewidth}{%
\begin{forest}
  chapter map,
  [Chapter 13
    [Directories [DASH protocol] [Sharing write back] [Owner-retaining variant]]
    [Token coherence [Counting tokens] [Owner token] [Persistent requests]]
    [Discussion [Organisation vs semantics] [Table C]]
  ]
\end{forest}}
\caption[]{Structure of Chapter~13. Both approaches change
how coherence messages travel and are ordered; the block-state rules on top
decide what happens to the dirty-sharing transaction.}
\label{fig:ch13-map}
\end{figure}

% ---------------------------------------------------------------
\section{Directory-Based Coherence}
\label{sec:scal-dir}
\sectionmap{[Directories [Replace broadcast by a sharer list] [DASH: three directory states]
  [Remote dirty read: four messages] [Memory is updated: sharing write back]]}
\index{MESI}\index{MOESI}\index{Owned state}\index{dirty sharing}\index{snooping protocol}\index{directory protocol}\index{DASH}\index{owner}%

Chapter~8 showed that snoop lookups grow as $N(N-1)$ (\cref{eq:snoop}) and that a
bus saturates once demand passes its capacity (\cref{eq:util}). A
\gls{directory}\index{directory protocol} replaces broadcast with a per-block record
of the caches that may hold the block, so requests reach only those
caches~\cite{censier1978directory}. A directory decides where a request goes and
where it is ordered. Whether a modified block may be shared is a separate
question, decided by the state protocol.

The Stanford DASH multiprocessor is the standard published example of a directory
protocol for distributed shared memory. Each memory block is uncached-remote,
shared-remote or dirty-remote in the directory, and each cached copy is invalid,
shared or dirty~\cite[\S4]{lenoski1990dash}. A directory entry has one state bit and
a bit vector with ``a bit for each of the sixteen clusters
supported''~\cite[\S3]{lenoski1990dash}, which by our count is 17 bits per memory
block. A request that reaches a cluster no longer holding the block gets a negative
acknowledgement and is sent again~\cite[\S4.1, \S4.6]{lenoski1990dash}.

The transaction of interest is a read of a block that is dirty in a remote cluster.
The home directory forwards the request ``to the owning, dirty cluster''. The owner
sends the data straight to the requester and also sends a sharing writeback to the
home: ``The sharing writeback request writes the cache block back to memory and
also updates the directory.'' The directory entry becomes shared-remote and the old
owner's copy becomes read-only~\cite[\S4.1]{lenoski1990dash}. The paper gives no
reason for updating memory at this point. This three-cluster transaction takes
about 80 processor clocks, against about 61 for a read served by a remote home with
no dirty copy~\cite[\S4.2, Fig.~6]{lenoski1990dash}. DASH thus performs, through
the home node and over a network, the same memory update that MESI performs on a
bus. Its cache states have no shared-dirty state, so a dirty block must become
clean when it is shared, and moving from snooping to a directory changes only the
route the transaction takes.

A directory can keep a dirty owner if it records which cache is the owner as well
as the sharers. A read of a dirty block is then forwarded to the owner, answered
from its cache and logged as a new sharer, with no writeback. Loughlin et al.\
study this combination on Intel's directory-based ccNUMA systems. They show that
MESI's ``downgrade writebacks'' on dirty sharing reach DRAM, that MOESI's \stO\
state removes them, and their MOESI-prime adds two states to avoid redundant
writes of the in-DRAM directory as well~\cite{loughlin2022moesiprime}.
\Cref{fig:dir-dirty-read} compares DASH with an owner-retaining directory. The
second is our construction, assuming one control message tells the home about the
new sharer.

\begin{figure}[htbp]
\centering
\begin{bigfig}[0.88]
\begin{tikzpicture}[font=\sffamily\scriptsize,
    ag/.style={draw, thick, rounded corners=2pt, minimum width=1.6cm,
               minimum height=0.55cm, font=\sffamily\small, align=center},
    m/.style={-{Stealth}, thick}]
  \node[font=\sffamily\bfseries\small] at (2.3,2.5) {(a) DASH: sharing write back};
  \node[ag, fill=mesiBlue!15, draw=mesiBlue] (R) at (0,0) {Req.};
  \node[ag, fill=moesiOrange!15, draw=moesiOrange] (H) at (2.3,1.6) {Home};
  \node[ag, fill=mesiBlue!15, draw=mesiBlue] (O) at (4.6,0) {Owner};
  \draw[m, coreGrey] (R) -- node[left, pos=0.4]{1. read 8\,B} (H);
  \draw[m, coreGrey] (H) -- node[left, pos=0.55]{2. fwd 8\,B} (O);
  \draw[m, mesiBlue] (O) -- node[below]{3. data 72\,B} (R);
  \draw[m, alertRed, very thick] (O.north east) to[out=70, in=0] node[right, text=alertRed, pos=0.4]{4. sharing WB 72\,B} (H.east);
  \node[text=coreGrey] at (2.3,-1.0) {4 messages, 160\,B; memory written};
  \begin{scope}[xshift=7.6cm]
  \node[font=\sffamily\bfseries\small] at (2.3,2.5) {(b) Directory with an owner};
  \node[ag, fill=mesiBlue!15, draw=mesiBlue] (R2) at (0,0) {Req.};
  \node[ag, fill=moesiOrange!15, draw=moesiOrange] (H2) at (2.3,1.6) {Home};
  \node[ag, fill=moesiOrange!15, draw=moesiOrange] (O2) at (4.6,0) {Owner (\stO)};
  \draw[m, coreGrey] (R2) -- node[left, pos=0.4]{1. read 8\,B} (H2);
  \draw[m, coreGrey] (H2) -- node[left, pos=0.55]{2. fwd 8\,B} (O2);
  \draw[m, moesiOrange] (O2) -- node[below]{3. data 72\,B} (R2);
  \draw[m, coreGrey, dashed] (O2.north east) to[out=70, in=0] node[right, pos=0.4]{4. ack 8\,B} (H2.east);
  \node[text=coreGrey] at (2.3,-1.0) {4 messages, 96\,B; memory not written};
  \end{scope}
\end{tikzpicture}
\end{bigfig}
\caption[A remote dirty read under two directory protocols]{A read of a block dirty
in another cache. (a)~DASH~\cite[\S4.1]{lenoski1990dash}: the sharing writeback
updates memory and the directory. (b)~An owner-retaining directory (our
construction): a control message records the new sharer and memory stays stale.
Both have three messages on the critical path; bytes use \cref{tab:msg-classes}.}
\label{fig:dir-dirty-read}
\end{figure}

% ---------------------------------------------------------------
\section{Token Coherence}
\label{sec:scal-token}
\sectionmap{[Token coherence [T tokens per block] [One token to read; all to write]
  [The owner token carries data] [Correctness substrate plus performance policy]]}
\index{MOESI}\index{directory protocol}\index{token coherence}\index{write-invalidate}\index{sharing patterns!migratory}\index{owner}\index{reading indices}%

Martin, Hill and Wood separate the mechanism that keeps a system correct from the
one that makes it fast~\cite{martin2003token}. The correctness substrate counts
tokens\index{token coherence}: ``the system assigns each block $T$ tokens (where $T$
is at least as large as the number of processors)'', and a processor ``can write a
block only if it holds all $T$ tokens'' and ``can read a block only if it holds at
least one token''~\cite[\S3.1]{martin2003token}. A writer must collect every token,
so nobody else can hold a readable copy during a write. This gives the
single-writer property of write-invalidate without an ordered interconnect.

One token is the owner token, and ``if a coherence message contains the owner
token, it must contain data''. Holding the owner token without all the others
``maps to the familiar MOESI state owned (O)''~\cite[\S3.1]{martin2003token}. On a
shared (read) request, a holder missing some tokens ``sends the data with one token
(usually not the owner token)''~\cite[\S4.2]{martin2003token}. On eviction a
processor returns all its tokens to memory, with data if the owner token is among
them~\cite[\S3.1]{martin2003token}. The paper does not state directly what happens
to memory on a read of a modified block. From these rules we infer that the owner
keeps the owner token and the duty to return data to memory, so the read is served
without a memory update, as in MOESI. One exception is stated: a processor that has
all tokens and has written the block answers a shared request with ``the data and
all tokens''~\cite[\S4.2]{martin2003token}, an optimisation for migratory sharing
that hands over ownership instead of sharing it.

Requests in the performance policy are ``fast, unordered `hint' requests'' that may
fail. A processor that suspects starvation issues a persistent request, and the
substrate activates ``at most one persistent request per
block''~\cite[\S3.2, \S4.1]{martin2003token}. The evaluated policy, TokenB,
broadcasts all transient requests over an unordered two-dimensional
torus~\cite[\S4.2, \S5.2]{martin2003token}. On a simulated 16-processor SPARC system
running OLTP, Apache and SPECjbb, ``only 3.0\,\% of cache misses are issued more
than once and only 0.2\,\% resort to persistent requests''. TokenB is reported
17--54\,\% faster than the paper's directory protocol and 8--29\,\% faster than its
Hammer protocol, while the directory protocol ``uses 21--25\,\% less'' interconnect
traffic~\cite[\S6]{martin2003token}. Token state costs $2+\log_2 T$ bits per block,
one byte for $T=64$ with 64-byte blocks~\cite[\S3.1]{martin2003token}.

TokenB removes the indirection of \cref{fig:dir-dirty-read}, since the owner answers
the requester directly, and by the inference above it also removes the writeback.
It pays with broadcast. Counting each broadcast as $N-1$ point-to-point messages of
8 bytes (the paper's request size~\cite[\S5.1]{martin2003token}), a read served by
an owner costs
\begin{equation}
  B_{\mathrm{TokenB}}(N) = 8\,(N-1) + 72 \ \text{bytes},
  \label{eq:tokenb}
\end{equation}
or 80, 96, 128 and 192 bytes for 2, 4, 8 and 16 processors. This is our analytical
count; the paper reports aggregate traffic, not this transaction.

% ---------------------------------------------------------------
\section{Limitations and Discussion}
\label{sec:scal-discussion}
\sectionmap{[Discussion [Directories cut broadcast and add a hop]
  [Tokens cut the hop and add broadcast] [Writeback set by the state rules] [Table C]]}
\index{MOESI}\index{directory protocol}\index{token coherence}\index{DASH}\index{owner}%

Directories and tokens answer the scaling problem of Chapter~8 in opposite ways. A
directory removes broadcast but adds an indirection and per-block storage that
grows with the processor count; Lenoski et al.\ note that full bit vectors grow with
both memory size and processor count and suggest limited-pointer and sparse
directories~\cite[\S6]{lenoski1990dash}. Token coherence removes the indirection and
needs no ordered interconnect, but uses broadcast requests and a starvation
mechanism. Neither, by its organisation alone, decides whether a modified block is
written back when shared. DASH writes it back, while an owner-retaining directory,
MOESI-prime and the owner token do not (\cref{tab:organisations}). Reducing
communication overhead and removing premature writebacks are therefore independent
goals, and published systems combine either kind of organisation with either kind
of state protocol.

\begin{table}[htbp]
\centering
\caption{Table C: coherence organisations, separating communication and
permissions from the dirty-sharing transaction.}
\label{tab:organisations}
\scriptsize
\begin{tabular}{@{}L{1.55cm} L{2.2cm} L{1.9cm} L{2.0cm} L{2.0cm} L{2.6cm}@{}}
\toprule
Organisation & Communication & Ordering point & Permission mechanism &
  Extra storage & Dirty-sharing transaction \\
\midrule
Snooping bus & broadcast to all caches & the bus & block states (MSI family)
  & none & set by the state protocol: MESI writes back, MOESI defers \\
Directory (DASH) & point-to-point via home & home directory & directory
  state + cache states & state bit + sharer vector per block~\cite{lenoski1990dash}
  & DASH writes back (sharing writeback)~\cite{lenoski1990dash}; an
  owner-retaining directory need not~\cite{loughlin2022moesiprime} \\
Token (TokenB) & broadcast hints, unordered torus & none needed; persistent
  requests for starvation & token counts and owner token & $2+\log_2T$ bits
  per block~\cite{martin2003token} & owner token keeps responsibility; no
  memory update inferred~\cite{martin2003token} \\
\bottomrule
\end{tabular}
\end{table}

% ===============================================================
%  Chapter 14 -- Comparative Analysis of Existing Solutions
% ===============================================================
\chapter{Comparative Analysis of Existing Solutions}
\label{ch:comparative}

\begin{figure}[htbp]
\centering
\adjustbox{max width=\linewidth}{%
\begin{forest}
  chapter map,
  [Chapter 14
    [Mechanisms [Three classes] [Nine approaches]]
    [Traffic [Counting rules] [Four sharing patterns]]
    [Trade-offs {\&} evaluations [{Latency, state}] [Seven studies]]
    [Research gap [Table D] [Two open questions]]
  ]
\end{forest}}
\caption[]{Structure of Chapter~14: one set of criteria, from
Chapter~10, applied to every approach of Chapters~11 to~13.}
\label{fig:ch14-map}
\end{figure}

% ---------------------------------------------------------------
\section{Comparison of Underlying Mechanisms}
\label{sec:cmp-mechanisms}
\sectionmap{[Mechanisms [Direct solutions change the states]
  [Alternative strategies change propagation] [Organisations change communication]
  [Nine approaches in one table]]}
\index{Dragon protocol}%

\Cref{fig:classes} places each approach in one of the three classes of Part~III.
Dragon appears twice because it combines an ownership state with update.
\Cref{tab:mechanisms} records, for each approach, the properties that decide its
behaviour on the transaction of \cref{fig:dirty-transaction}. The last two rows are
organisations, not complete state protocols, so their entries describe the
published protocols evaluated with them.

\begin{figure}[htbp]
\centering
\adjustbox{max width=0.82\linewidth}{%
\begin{forest}
  for tree={
    grow'=0, parent anchor=east, child anchor=west, anchor=west,
    align=center, font=\sffamily\small, draw, thick, rounded corners=3pt,
    inner sep=4pt, l sep+=14pt, s sep+=1pt,
    edge={thick, coreGrey!60, -{Stealth[length=5pt]}},
  },
  [Approaches\\reviewed, fill=mesiBlue, draw=mesiBlue, text=white, font=\sffamily\bfseries\small
    [Direct solutions:\\shared-dirty owner, fill=memGreen!20, draw=memGreen
      [Berkeley, fill=memGreen!8, draw=memGreen!60]
      [MOESI class, fill=memGreen!8, draw=memGreen!60]
      [Dragon (owner part), fill=memGreen!8, draw=memGreen!60]]
    [Alternative strategies:\\propagation and forwarding, fill=moesiOrange!20, draw=moesiOrange
      [Firefly: update and write-through, fill=moesiOrange!8, draw=moesiOrange!60]
      [Dragon (update part), fill=moesiOrange!8, draw=moesiOrange!60]
      [MESIF: clean forwarding, fill=moesiOrange!8, draw=moesiOrange!60]]
    [Organisations:\\how requests travel, fill=coreGrey!20, draw=coreGrey
      [Directory (DASH), fill=coreGrey!8, draw=coreGrey!60]
      [Token coherence, fill=coreGrey!8, draw=coreGrey!60]]
  ]
\end{forest}}
\caption[Classification of the approaches reviewed]{The three classes. Only the
first targets the forced writeback by design; the second affects it as a side
effect; the third leaves it to the state protocol on top.}
\label{fig:classes}
\end{figure}

\begin{table}[htbp]
\centering
\caption{Underlying mechanisms of the approaches reviewed. ``WB'' = writeback.}
\label{tab:mechanisms}
\scriptsize
\begin{tabular}{@{}L{1.25cm} L{1.35cm} L{1.75cm} L{1.05cm} L{1.05cm} L{1.75cm} L{1.85cm} L{1.75cm}@{}}
\toprule
Approach & Strategy & State or permission model & Dirty data shared? &
  Explicit owner? & Supplier on remote read & Memory updated & Forced WB \\
\midrule
MSI & invalidate & M, S, I~\cite{sarangi2025basic} & no & M only & holder of M &
  at M$\to$S and on eviction of M & present \\
MESI & invalidate & M, E, S, I~\cite{papamarcos1984illinois} & no & M only & holder of M &
  at M$\to$S and on eviction of M & present (baseline) \\
MOESI & invalidate & M, O, E, S, I~\cite{sweazey1986moesi} & yes (O) & yes &
  owner (M or O) & on eviction of owner & eliminated; deferred to eviction \\
Berkeley & invalidate & I, Valid, SD, Dirty~\cite{archibald1986evaluation} & yes (SD)
  & yes & owner & on eviction of owner & eliminated; deferred to eviction \\
Dragon & update & VE, SC, SD, Dirty~\cite{archibald1986evaluation} & yes (SD)
  & yes & owner & on eviction of SD owner & eliminated; replaced by update traffic \\
Firefly & update, write-through & Dirty, Shared bits~\cite{thacker1987firefly} &
  no & no & caches holding it & at sharing, on every shared write, on eviction & retained; extra memory writes \\
MESIF & invalidate, forwarding & M, E, S, I, F~\cite{goodman2009mesif} & no (basic)
  & no & F holder (clean); M holder by migration & not stated for the M case &
  avoided by migration (basic) \\
Directory (DASH) & organisation & 3 directory + 3 cache states~\cite{lenoski1990dash}
  & no & dirty owner only & owner, forwarded by home & sharing WB at the remote read &
  retained \\
Token (TokenB) & organisation & token counts, owner token~\cite{martin2003token}
  & yes (owner token) & yes & owner-token holder & on eviction of owner token &
  eliminated (inferred) \\
\bottomrule
\end{tabular}
\end{table}

% ---------------------------------------------------------------
\section{Comparison of Communication and Memory Traffic}
\label{sec:cmp-traffic}
\sectionmap{[Traffic [One set of counting rules] [Bus protocols in one table]
  [Point-to-point counted separately] [Bar graphs per pattern]]}
\index{MESI}\index{MOESI}\index{MESIF}\index{Exclusive state}\index{directory protocol}\index{token coherence}\index{Dragon protocol}\index{sharing patterns!producer--consumer}\index{sharing patterns!migratory}\index{owner}%

None of the sources reports traffic for these protocols in one common
configuration, so this section is our own analytical accounting. It applies the
rules of \cref{tab:counting-rules} to the same access sequences under each
protocol. Bus and point-to-point protocols are counted separately, since one bus
broadcast and $N-1$ point-to-point messages are different costs.

\begin{table}[htbp]
\centering
\caption{Counting rules for the analytical comparison. Rows~1--3 are those of
\cref{tab:msg-classes}; rows~4--6 are assumptions added for Part~III.}
\label{tab:counting-rules}
\small
\begin{tabular}{@{}c L{5.4cm} L{6.0cm}@{}}
\toprule
 & Item & Rule \\
\midrule
1 & Control message (request, invalidation, upgrade) & 8\,B on the interconnect \\
2 & Data message (block to a cache) & 72\,B (64\,B block + 8\,B header) \\
3 & Writeback message & 72\,B on the interconnect; 64\,B written to memory \\
4 & Word update (Dragon, Firefly) & 16\,B (8\,B header + 8\,B word); Firefly
  also writes 8\,B to memory \\
5 & Bus broadcast & counted once: one transaction reaches all caches \\
6 & Point-to-point broadcast (TokenB) & counted as $N-1$ messages \\
\bottomrule
\end{tabular}
\end{table}

\Cref{tab:traffic-patterns} gives the cost of each sharing pattern of
\cref{tab:patterns} under the six bus protocols, as interconnect bytes and bytes
written to memory, derived from the transitions cited in Chapters~11 and~12. A
migratory hand-off under MESI, for example, is a read miss served by the \stM\
holder with a writeback ($8+72+72$) followed by an upgrade (8), 160\,B in all.
Under MOESI the owner serves the read without a writeback ($8+72$) and the upgrade
invalidates the old owner, whose data is superseded (8), giving 88\,B. Under
Dragon every core keeps a copy once it has taken part, so in steady state a
hand-off is one update; the writer becomes \textsc{Shared-Dirty} and the previous
owner \textsc{Shared-Clean}~\cite[p.~282]{archibald1986evaluation}.

\begin{table}[htbp]
\centering
\caption{Analytical traffic per unit of each sharing pattern, as interconnect
bytes\,/\,bytes written to memory, with one producer write per round ($w=1$) and
one consumer ($C=1$). Our counts, using \cref{tab:counting-rules}.}
\label{tab:traffic-patterns}
\small
\begin{tabular}{@{}L{3.0cm} c c c c c c@{}}
\toprule
Pattern (unit) & MSI & MESI & MOESI & Berkeley & Dragon & Firefly \\
\midrule
Private (read, write, evict) & 160\,/\,64 & 152\,/\,64 & 152\,/\,64 & 160\,/\,64 & 152\,/\,64 & 152\,/\,64 \\
Read-only (per reader) & 80\,/\,0 & 80\,/\,0 & 80\,/\,0 & 80\,/\,0 & 80\,/\,0 & 80\,/\,0 \\
Producer--consumer (per round) & 160\,/\,64 & 160\,/\,64 & 88\,/\,0$^{*}$ & 88\,/\,0$^{*}$ & 16\,/\,0$^{*}$ & 16\,/\,8 \\
Migratory (per hand-off) & 160\,/\,64 & 160\,/\,64 & 88\,/\,0$^{*}$ & 88\,/\,0$^{*}$ & 16\,/\,0$^{*}$ & 16\,/\,8 \\
\bottomrule
\multicolumn{7}{@{}p{12.6cm}@{}}{\footnotesize $^{*}$Plus one deferred
writeback (72\,B, 64\,B to memory) when the owner finally evicts the block,
amortised over all rounds. Private-pattern differences come only from the
\stE\ state: MSI and Berkeley need an upgrade on the first write; the others
upgrade silently~\cite[pp.~279--282]{archibald1986evaluation}.}
\end{tabular}
\end{table}

The private and read-only patterns cost the same under every protocol apart from
the \stE\ effect, so a workload made of them gains nothing from ownership or
update. In producer--consumer and migratory sharing, ownership removes every
memory write and cuts interconnect bytes by 45\,\% relative to MESI, which is the
forced-writeback fraction of \cref{eq:pc-model}. Update protocols move the fewest
bytes at $w=1$, but lose that advantage as $w$ grows
(\cref{fig:update-vs-invalidate}). \Cref{fig:bars-pc,fig:bars-mig} show the two
patterns in which the protocols differ.

\begin{figure}[htbp]
\centering
\begin{minipage}[t]{0.5\linewidth}\centering
\begin{tikzpicture}
  \begin{axis}[width=\linewidth, height=5.2cm, ybar, bar width=4pt,
      xlabel={Number of consumers $C$}, ylabel={Bytes per round},
      symbolic x coords={1,2,4,8}, xtick=data, ymin=0, ymax=800,
      enlarge x limits=0.18,
      legend style={at={(0.02,0.98)}, anchor=north west, font=\tiny},
      grid=major, grid style={coreGrey!15},
      label style={font=\footnotesize}, tick label style={font=\scriptsize}]
    \addplot[fill=mesiBlue!65, draw=mesiBlue] coordinates {(1,160) (2,240) (4,400) (8,720)};
    \addlegendentry{MSI / MESI}
    \addplot[fill=moesiOrange!65, draw=moesiOrange] coordinates {(1,88) (2,168) (4,328) (8,648)};
    \addlegendentry{MOESI / Berkeley}
    \addplot[fill=tealC!70, draw=tealC] coordinates {(1,16) (2,16) (4,16) (8,16)};
    \addlegendentry{Dragon / Firefly, $w=1$}
    \addplot[fill=plumC!45, draw=plumC] coordinates {(1,128) (2,128) (4,128) (8,128)};
    \addlegendentry{Dragon / Firefly, $w=8$}
  \end{axis}
\end{tikzpicture}
\caption[Producer--consumer traffic by protocol]{Producer--consumer: ownership saves
72 bytes per round; update cost does not depend on $C$ but grows with $w$. Firefly
also writes $8w$ bytes to memory. Our model.}
\label{fig:bars-pc}
\end{minipage}\hfill
\begin{minipage}[t]{0.48\linewidth}\centering
\begin{tikzpicture}
  \begin{axis}[width=\linewidth, height=5.2cm, ybar, bar width=4pt,
      ylabel={Bytes per hand-off},
      symbolic x coords={MSI,MESI,MOESI,Berk.,Dragon,Firefly}, xtick=data,
      ymin=0, ymax=220, enlarge x limits=0.1,
      legend style={at={(0.98,0.98)}, anchor=north east, font=\tiny},
      grid=major, grid style={coreGrey!15},
      label style={font=\footnotesize}, tick label style={font=\tiny}]
    \addplot[fill=skyC!60, draw=skyC] coordinates
      {(MSI,160) (MESI,160) (MOESI,88) (Berk.,88) (Dragon,16) (Firefly,16)};
    \addlegendentry{interconnect, $w=1$}
    \addplot[fill=amberC!45, draw=amberC] coordinates
      {(MSI,160) (MESI,160) (MOESI,88) (Berk.,88) (Dragon,128) (Firefly,128)};
    \addlegendentry{interconnect, $w=8$}
    \addplot[fill=brickC!65, draw=brickC] coordinates
      {(MSI,64) (MESI,64) (MOESI,0) (Berk.,0) (Dragon,0) (Firefly,8)};
    \addlegendentry{to memory, $w=1$}
  \end{axis}
\end{tikzpicture}
\caption[Migratory-sharing traffic by protocol]{Migratory: bytes per hand-off when
each core reads a block and writes it $w$ times. Only MSI and MESI write memory at
every hand-off. Our model.}
\label{fig:bars-mig}
\end{minipage}
\end{figure}

For the point-to-point organisations of Chapter~13 the dirty-sharing transaction is
compared in \cref{tab:p2p}. Removing the writeback saves the same 64 bytes of memory
traffic in the directory and token organisations, while TokenB's interconnect cost
grows with $N$ and the directory's does not. We do not cost MESIF: its message set
(invalid-state acknowledgements, cancellations and data acknowledgements to the
home node~\cite[\S2.1.1]{goodman2009mesif}) belongs to a protocol the report gives
only in outline, and a byte count would rest on assumptions the source does not
support.

\begin{table}[htbp]
\centering
\caption{The remote dirty read on point-to-point organisations. Our analytical
counts, using \cref{tab:counting-rules}.}
\label{tab:p2p}
\small
\begin{tabular}{@{}L{3.9cm} c c c c@{}}
\toprule
Organisation & Messages & Bytes & Critical-path hops & Memory written \\
\midrule
Directory, DASH (sharing WB)~\cite{lenoski1990dash} & 4 & 160 & 3 & 64\,B \\
Directory with owner (\cref{fig:dir-dirty-read}b) & 4 & 96 & 3 & 0 \\
TokenB, $N=4$ (\cref{eq:tokenb}) & 4 & 96 & 2 & 0 \\
TokenB, $N=16$ (\cref{eq:tokenb}) & 16 & 192 & 2 & 0 \\
\bottomrule
\end{tabular}
\end{table}

% ---------------------------------------------------------------
\section{Trade-offs and Published Evaluations}
\label{sec:cmp-tradeoffs}
\label{sec:cmp-evaluations}
\sectionmap{[Trade-offs [Latency evidence from the papers] [State and storage cost]
  [Fewer write backs need not mean less time] [Seven studies]]}
\index{MESI}\index{MOESI}\index{MESIF}\index{Exclusive state}\index{forced writeback}\index{token coherence}\index{Berkeley protocol}\index{Dragon protocol}\index{Firefly protocol}\index{Illinois protocol}\index{Futurebus}\index{DASH}\index{gem5}\index{cache-to-cache transfer}%

The sources report latency in different units and configurations, so only
comparisons within one study mean anything. On Futurebus a protocol that must
update memory during a cache-to-cache transfer aborts and restarts the
transaction~\cite[p.~418]{sweazey1986moesi}, so the forced writeback lengthens the
read itself. In DASH a read involving a dirty remote cluster takes about 80
processor clocks against about 61 for a clean remote
read~\cite[\S4.2, Fig.~6]{lenoski1990dash}. MESIF cuts a read of a line held in E, M
or F to two hops but offers ``no benefits over the 3-hop protocols'' for lines in I
or S~\cite[\S3]{goodman2009mesif}. TokenB answers in two hops but reissues 3.0\,\%
of misses~\cite[\S6]{martin2003token}. The state and storage costs are in row~10 of
\cref{tab:matrix}: three bits per line plus the \textsc{DI} and \textsc{CH} signals
for MOESI~\cite[p.~416]{sweazey1986moesi}, a 17-bit entry per block for 16-cluster
DASH~\cite[\S3]{lenoski1990dash}, and $2+\log_2T$ bits plus a persistent-request table of 8 bytes per arbiter
for tokens~\cite[\S3.1--3.2]{martin2003token}.

Removing a memory writeback does not by itself reduce total traffic or execution
time. The saved bytes can be replaced by update traffic (Dragon), kept alongside it
(Firefly) or traded for broadcast requests (TokenB), and saved traffic becomes
saved time only when the interconnect or memory was the bottleneck
(\cref{fig:bus-util}). MOESI-prime is the clearest published case: it removed MESI's
downgrade writebacks and reached its goal of fewer DRAM row activations, while
execution time changed by between $-0.51\,\%$ and
$+0.61\,\%$~\cite{loughlin2022moesiprime}.

\Cref{tab:evaluations} lists the comparative studies we found. Their set-ups range
from a synthetic single-bus model and queueing analysis to full-system simulation of
commercial workloads and gem5 runs with different protocol implementations, so
their numbers cannot be compared with each other and we do not compare them.
Stenstr{\"o}m~\cite{stenstrom1990survey} surveys hardware and software coherence
schemes of the same period. Two points matter for this report. The only early study
that compared an ownership protocol with an invalidate protocol under identical
conditions, Archibald and Baer, compared whole protocols: Berkeley and Illinois
differ both in having \stE\ and in ownership (\cref{tab:berkeley-states}), so the
effect of ownership alone cannot be read from their
results~\cite{archibald1986evaluation}. And the recent gem5 studies compare
implementations that differ in many ways; in the case reported by Mallya et al.\
the MOESI implementation did not share dirty data on our
transaction~\cite{mallya2015simulation}.

\begin{table}[htbp]
\centering
\caption{Published comparative evaluations relevant to the forced writeback.}
\label{tab:evaluations}
\scriptsize
\begin{tabular}{@{}L{1.75cm} L{2.0cm} L{2.0cm} L{2.15cm} L{2.6cm} L{2.0cm}@{}}
\toprule
Study & Protocols & Method and system & Workload, metric & Principal finding &
  Isolates the O state? \\
\midrule
Archibald {\&} Baer 1986~\cite{archibald1986evaluation} & Write-once, Synapse,
  Berkeley, Illinois, Firefly, Dragon & simulation; one bus; 1--15 processors &
  synthetic streams; ``system power'' & update best for shared data; Berkeley
  beats Illinois at high sharing (p.~295) & no: whole protocols differ \\
Yang et al.\ 1989~\cite{yang1989analysis} & the six above and write-through &
  queueing model (MVA); packet-switched bus & analytical & full text not
  consulted & --- \\
Martin et al.\ 2003~\cite{martin2003token} & TokenB, snooping, directory,
  Hammer & Simics full system; 16 processors & OLTP, Apache, SPECjbb; runtime,
  traffic & TokenB faster; directory 21--25\,\% less traffic & no: organisations
  compared \\
Goodman {\&} Hum 2009~\cite{goodman2009mesif} & MESIF vs.\ 3-hop & in-house MVA,
  undisclosed & projection & 6--11\,\% (4 nodes) & no \\
Mallya et al.\ 2015~\cite{mallya2015simulation} & MI, MESI, MOESI & gem5 Ruby;
  2--4 nodes & SPLASH-2; misses, energy & MESI better in some cases & no:
  different gem5 protocols \\
Loughlin et al.\ 2022~\cite{loughlin2022moesiprime} & MESI, MOESI, MOESI-prime &
  gem5; 2--8 nodes & PARSEC, SPLASH-2x; row activations & MOESI removes downgrade
  writebacks; runtime within $\pm0.61\,\%$ & not established \\
Fagundes {\&} Souza 2024~\cite{fagundes2024evaluation} & MESI, MOESI & gem5 full
  system; 4, 8 cores & one OpenMP kernel ($\pi$) & MOESI better, less so with
  larger L2 & no \\
\bottomrule
\end{tabular}
\end{table}

% ---------------------------------------------------------------
\section{Research Gap and Motivation for Our Contribution}
\label{sec:cmp-gap}
\sectionmap{[Gap [What the literature establishes] [Table D]
  [Open 1: isolate the O state] [Open 2: predict the benefit]]}
\index{MESI}\index{MOESI}\index{Owned state}\index{sharing patterns!migratory}\index{gem5}%

The sources establish four points. Ownership protocols let modified data be shared
without updating memory and stay coherent~\cite{sweazey1986moesi,archibald1986evaluation},
so the timing of the memory update is a design choice. Different approaches trade
memory traffic, interconnect traffic and hardware state against one another
(\cref{tab:traffic-patterns,tab:p2p}). And the benefit of each depends on the
workload's sharing
behaviour~\cite[p.~295]{archibald1986evaluation}\cite{fagundes2024evaluation,loughlin2022moesiprime}.
\Cref{tab:matrix} judges every approach against the requirements of Chapter~10.

\begin{table}[htbp]
\centering
\caption{Table D: every approach against the requirements of Chapter~10.
DA: directly addresses; PA: partially addresses; RP: addresses a related
problem; OR: orthogonal; NE: not established by the available source.}
\label{tab:matrix}
\scriptsize
\begin{tabular}{@{}L{2.05cm} L{1.4cm} L{1.4cm} L{1.4cm} L{1.4cm} L{1.4cm} L{1.4cm} L{1.4cm}@{}}
\toprule
Criterion & MOESI & Berkeley & Dragon & Firefly & MESIF & Directory (DASH) & Token \\
\midrule
1 Completion & yes$^{a}$ & yes$^{a}$ & yes$^{a}$ & yes$^{a}$ & yes$^{a}$ & yes$^{a}$ & yes$^{a}$ \\
2 Order & yes$^{a}$ & yes$^{a}$ & yes$^{a}$ & yes$^{a}$ & yes$^{a}$ & yes$^{a}$ & yes$^{a}$ \\
3 Removes the immediate forced WB & DA: deferred & DA: deferred & DA: deferred &
  no: retained & PA: by migration & no: sharing WB & DA (inferred) \\
4 Keeps silent E$\to$M & yes & no: no E & yes (VE) & yes (VE) & E present; NE &
  OR: no E state & NE \\
5 Shares modified data & yes & yes & yes & no & no; FM proposed & no & yes \\
6 Explicit ownership & yes & yes & yes & no & no & dirty owner only & yes \\
7 Extra control traffic & none over MESI & invalidation on first write &
  update per shared write & update per shared write & home acks (IACK, CNCL,
  DACK) & forward + acks & broadcast; 3\,\% reissue \\
8 Extra data traffic & none & none & none & none & producer re-fetch (our
  analysis) & sharing WB & none (inferred) \\
9 Memory bandwidth & lower & lower & lower & higher & NE & as MESI & lower
  (inferred) \\
10 Hardware & 3 bits; DI, CH & 2 bits; interlocks & update path & update path;
  MShared & 3 bits; 3 VCs & 17 bits per block & $2+\log_2T$ bits \\
11 Scalability & bus-based as published & bus-based & bus-based & bus-based &
  point-to-point & scalable & unordered; broadcast \\
12 Workload sensitivity & high-sharing gains & high-sharing gains & writes per
  read & writes per read & E, M, F hits only & NE & commercial only \\
\bottomrule
\multicolumn{8}{@{}p{12.6cm}@{}}{\footnotesize $^{a}$As published; each
protocol is presented by its authors as coherent. Formal verification is
outside the scope of this report (\cref{sec:prob-scope}).}
\end{tabular}
\end{table}

Two questions remain open, and Part~IV takes them up.

\paragraph{Isolating the effect of the Owned state.} None of the studies in
\cref{tab:evaluations} compares MESI and MOESI with \stO\ as the only difference.
The early comparison varied whole protocols, and the gem5 comparisons used
implementations that differ in many states and transitions, in one case without
the \stM$\to$\stO\ transition at all. We do not claim that no such study exists
anywhere, only that the studies we found do not provide one, so their results
cannot be credited to the \stO\ state alone. Part~IV therefore builds a controlled
simulator in which MESI and MOESI share every component and setting, leaving the
\stM$\to$\stO\ behaviour as the single changed variable.

\paragraph{Predicting the benefit from MESI behaviour.} The sources agree that the
benefit of ownership depends on sharing, and Archibald and Baer parameterise their
model by the fraction of references to shared
blocks~\cite[p.~286]{archibald1986evaluation}. Detecting sharing patterns is also
known: adaptive protocols classify blocks as migratory at run time and change how
they handle them~\cite{stenstrom1993adaptive}, using rules such as spotting a write
by a processor other than the last writer~\cite{nilsson1994adaptive}. We did not
find a workload measure, taken under MESI, that has been proposed and tested as a
predictor of the gain from moving to MOESI. Part~IV defines the Dirty Sharing
Ratio, the fraction of coherence transactions in which a read is served by a cache
holding the block in \stM, for this purpose. It differs from adaptive detection in
being an offline characterisation used to choose a protocol, rather than a run-time
mechanism inside one.
